\section{Ограничения и бъдещо развитие}
\label{sec:future_work}

Текущата версия на библиотеката вече демонстрира работеща архитектура за библиотека за обектно-релационно съпоставяне по време на компилация върху релационни и нерелационни свързващи компоненти, но научната коректност изисква ясно разграничение между вече постигнатото и следващите стъпки. Библиотеката съдържа реален слой за междинно представяне на заявките, базиран на възможности конекторен договор, миграционен прототип, помощници (помощници) за нишкова безопасност и набор от модулни и интеграционни тестове. Въпреки това част от подсистемите все още трябва да бъдат разширени или валидирани по-задълбочено в условия, близки до производствените (условия, близки до производствените).

Бъдещото развитие не е произволен списък от идеи. То следва непосредствено от архитектурните решения, анализирани в предходните глави. Ако C++ моделът е първичен носител на логическата схема, следваща естествена стъпка е по-пълна автоматизация на инструментариума, съобразен със схемата (инструментариум, съобразен със схемата). Ако конекторният слой е централна точка на разширяемост (разширяемост), тогава бъдещата работа трябва да уточни по-фини класове от възможности. Ако междинното представяне е реално споделен слой между различни свързващи компоненти, тогава то трябва да бъде изпитано при по-богати експлоатационни сценарии и ограничения на производителността.

\subsection{Стратегически направления за развитие}

От гледна точка на настоящата дипломна работа бъдещото развитие може да бъде групирано в четири стратегически направления. Първото е \emph{по-голяма експлоатационна зрялост} на вече наличните целеви системи. Второто е \emph{по-пълна автоматизация на схемата и метаданните}. Третото е \emph{намаляване на режийните разходи по време на изпълнение (по време на изпълнение режийни разходи)} чрез по-ниско ниво на изпълнение и по-добър контрол върху слоя на мрежовия протокол (мрежов протокол). Четвъртото е \emph{намаляване на шаблонния код (boilerplate)} чрез по-дълбока интеграция с новите езикови механизми на C++.

\begin{figure}[H]
\centering
\begin{tikzpicture}[node distance=1.35cm and 1.2cm]
    \node[ormaccent] (core) {\textbf{Текуща библиотека}\\типизирано междинно представяне + договор на конектора + тестове};
    \node[ormblock, above right=of core] (maturity) {Експлоатационна\\зрялост};
    \node[ormblock, below right=of core] (schema) {Автоматизация,\\съобразена със схемата};
    \node[ormblock, below left=of core] (perf) {Оптимизация на\\изпълнението / мрежовия протокол};
    \node[ormblock, above left=of core] (lang) {Автоматизация\\на езиково ниво};

    \draw[ormline] (core) -- (maturity);
    \draw[ormline] (core) -- (schema);
    \draw[ormline] (core) -- (perf);
    \draw[ormline] (core) -- (lang);
\end{tikzpicture}
\caption{Четири стратегически направления за бъдещо развитие}
\label{fig:future_work_axes}
\end{figure}

Фигура~\ref{fig:future_work_axes} е полезна, защото показва, че бъдещата работа не е еднопластова. Част от задачите са насочени към експлоатационна зрялост, други --- към по-добра автоматизация по време на компилация, а трети --- към инженеринг на производителността и времето за изпълнение. Това разграничение е важно, защото различните направления изискват различни критерии за успех.

\subsection{Разширяване на покритието на целеви системи на живо}

Макар хранилището вече да съдържа интеграционни тестове на живо за няколко системи, зрелостта на отделните реализации не е еднаква. SQLite остава най-естественият референтен релационен път, защото съчетава най-плътно свързани рендиране на заявки, логика за свързване (свързване) и попълване на резултати (попълване на резултати). При PostgreSQL, MySQL, MongoDB, Redis, Neo4j и Cassandra следващата стъпка е по-широко покриване на експлоатационни сценарии: по-богати типове, диагностични пътища при грешка, политики за жизнен цикъл на връзките и оценка на поведението при по-сложни заявки.

Това направление е важно не само като инженерно подобрение, а и като научна проверка на границите на общото междинно представяне. Само при по-широко емпирично покритие може да се оцени доколко един и същ логически слой остава адекватен за различни типове системи. Например логика, базирана на съединявания, която е естествена за PostgreSQL, има съвсем различен експлоатационен смисъл в Neo4j или Cassandra. Следователно покритието на живо не е просто тестова дисциплина, а средство за изследване на архитектурната приложимост.

Допълнително, покритието на живо трябва да обхване и негативните пътища: грешки при установяване на връзка, частично невалидни параметри, специфични за драйвера сценарии за изтичане на времето (изтичане на времето) и поведение при прекъсване на връзката. Тези въпроси почти винаги са подценени в ранните ОРС библиотеки, но в контекст с множество свързващи компоненти те са решаващи за реалната употреба.

\subsection{Пълна интеграция, съобразена със схемата, и интроспекция по време на изпълнение}

Прототипният миграционен слой е достатъчен, за да покаже как C++ моделът на същността може да участва в описване на еволюцията на схемата. Следващото естествено направление е интроспекция на живо на реалната схема, извличане на метаданни от драйверите и безопасно изпълнение на DDL операции. Това е особено важно за релационните целеви системи, но не е без значение и за документни или широки колонни системи, където понятието за схема може да бъде по-гъвкаво или частично.

По-нататъшната работа в тази посока трябва да изясни и границата между гаранции по време на компилация и факти по време на изпълнение за вече разположена система. Компилаторът може да валидира структурата на заявката, но не и реалното състояние на производствената схема без допълнителна информация по време на изпълнение. Следователно бъдещият миграционен слой трябва да комбинира два източника на истина: C++ описанието на същността и извлечената картина от метаданни на реалната база.

\begin{figure}[H]
\centering
\begin{tikzpicture}[node distance=1.3cm and 1.1cm]
    \node[ormaccent] (model) {C++ модел\\метаданни на същности};
    \node[ormblock, right=of model] (live) {Целева система на живо\\метаданни};
    \node[ormblock, below=of $(model)!0.5!(live)$] (compare) {Разлика в схемата +\\анализ на съвместимост};
    \node[ormblock, below left=of compare] (ddl) {Генериране на DDL};
    \node[ormblock, below right=of compare] (report) {Доклад за безопасност\\и пробно изпълнение};
    \node[ormblock, below=of $(ddl)!0.5!(report)$] (apply) {Контролирано изпълнение\\на миграцията};

    \draw[ormline] (model) -- (compare);
    \draw[ormline] (live) -- (compare);
    \draw[ormline] (compare) -- (ddl);
    \draw[ormline] (compare) -- (report);
    \draw[ormline] (ddl) -- (apply);
    \draw[ormline] (report) -- (apply);
\end{tikzpicture}
\caption{Целева посока за инструментариум, съобразен със схемата, отвъд текущия миграционен прототип}
\label{fig:future_schema_tooling}
\end{figure}

Фигура~\ref{fig:future_schema_tooling} показва, че бъдещият миграционен конвейер трябва да бъде двупосочен: едновременно да чете от описанието по време на компилация и от живата система. Това е естествена следваща стъпка, ако библиотеката трябва да бъде използвана не само експериментално, но и като практичен инструмент за разработка.

\subsection{Нишкова безопасност и жизнен цикъл на връзките}

Съществуващите помощници (помощници) за пул от връзки, достъп, локален за нишката, и предпазител на транзакция показват, че библиотеката вече съдържа архитектурна основа за многонишкова употреба. Бъдещото развитие трябва да отговори на по-конкретни въпроси: как се дефинира \code{supports\_concurrent\_execute} за различните драйвери; кога една връзка е безопасна за споделяне; какви са правилата за повторно свързване (повторно свързване), затваряне и отложено освобождаване на ресурси.

Тук има и важен изследователски аспект: различните целеви системи имат различна семантика на конкурентността, а следователно обозначаването на възможности по време на компилация трябва да бъде достатъчно прецизно, за да не създава подвеждащи обещания. Напълно възможно е бъдещи версии на библиотеката да разделят общото \code{supports\_concurrent\_execute} на по-фини класове от възможности като безопасно паралелно четене, паралелно изпълнение с отделни манипулатори на връзки или строга политика за сериализация.

Практически това означава, че помощният слой за нишкова безопасност не трябва да остане изолиран модул. Напротив, той трябва да се превърне в архитектурно продължение на договора на конектора и да участва в пълната оценка на зрелостта на дадена целева система.

\subsection{Ниско ниво на изпълнение и оптимизации на мрежовия протокол}

Експерименталният слой в \path{WireProtocol/wire_protocol.hpp} показва възможни бъдещи посоки като изчакваеми обекти за \code{io\_uring} / IOCP, \code{zero\_copy\_result}, \code{batch\_insert} и \code{make\_constexpr\_sql}. Тези елементи все още не са готов за производство комуникационен стек, но очертават възможност библиотеката да се разшири отвъд класическия синхронен драйверен модел.

Особено перспективни са нулево-копийната обработка на резултати и генерирането на SQL по време на компилация за конектори, които го позволяват. Те биха приближили библиотеката до още по-ниски режийни разходи по време на изпълнение, без да се жертва типизираната архитектура. Това е важно, защото една честа критика към богато типизираните слоеве на абстракция е, че скриват прекомерен разход по време на изпълнение. Бъдещата работа трябва да покаже, че изразителността по време на компилация и ниските режийни разходи не са задължително противоположни цели.

Освен това темата за производителността не трябва да се разглежда само през латентност или пропускателна способност. Тя включва и разположението в паметта на резултатите, броя временни заделяния на памет (заделяния на памет), повторната употреба на артефакти на подготвени заявки и качеството на стратегиите за масово обработване (масово обработване) на ниво драйвер.

\subsection{Интеграция с отражение (отражение) в C++26 и допълнителна автоматизация}

Поддръжката на отражение в C++26 е вече концептуално подготвена чрез логика за откриване (откриване) и помощници, а \path{tests/модулни/test_cpp26_отражение.cpp} показва наличието на ранен път за верификация. Публичният програмен интерфейс обаче все още използва експлицитни низови буквални имена за \code{property<T, Name>}. Бъдещото развитие трябва да превърне пътя с отражение в стабилен публичен механизъм, който позволява автоматично извличане на имената на полетата, без това да нарушава обратната съвместимост.

Подобно развитие би намалило шаблонния код и би направило още по-видима идеята, че C++ моделът е първичният източник на логическата схема. От изследователска гледна точка то би било интересно и защото би доближило библиотеката до стил, ориентиран към домейна (ориентиран към модела), при който голяма част от информацията за метаданните се извлича автоматично, а не се описва повторно.

Тук обаче има и важна граница. Отражението не трябва да разрушава яснотата на договора по време на компилация. Ако автоматизацията намали прозрачността на типа или затрудни съобщенията за грешка, тя може да отслаби една от основните ценности на библиотеката. Следователно тази посока трябва да се развива внимателно.

\subsection{Разширяване на договора на конектора и таксономията на възможностите}

С добавянето на нови целеви системи ще стане необходимо по-прецизно разграничаване между различни класове възможности. Възможно е например бъдещи версии на библиотеката да разграничават отделно възможност за обхождане на графи, поддръжка на частично вграждане на документи, поддръжка на интроспекция на схемата, поддръжка на поточни резултати или изрична поддръжка на масово обработване (масово обработване). Тази посока следва естествено от текущата архитектура, базирана на характеристики.

По-нататъшната работа трябва да пази едно важно правило: новите възможности не бива да бъдат декоративни. Те трябва да носят реално значение по време на компилация и да ограничават недопустимите операции. Ако label за възможност няма архитектурно последствие, той не добавя стойност. Точно обратното --- прави договора по-шумен и по-малко точен.

\begin{table}[H]
\centering
\small
\caption{Приоритетни направления за бъдещо развитие}
\label{tab:future_work_priorities}
\begin{tabularx}{\textwidth}{L{3.0cm}L{3.2cm}L{3.1cm}Y}
\toprule
\textbf{Направление} & \textbf{Засегнати подсистеми} & \textbf{Очакван ефект} & \textbf{Основно предизвикателство} \\
\midrule
Зрялост на компонентите на живо & конектори, интеграционни тестове, пътища за грешки & по-висока практическа надеждност & различна експлоатационна семантика между компонентите \\
Автоматизация, съобразена със схемата & \code{migrate<DB>}, DDL куки, извличане на метаданни & по-близка до производствената ОРС интеграция & съчетаване на истина по време на компилация и по време на изпълнение \\
Оптимизация на изпълнението & слой на мрежовия протокол, подготвени заявки, масово обработване & по-ниски режийни разходи по време на изпълнение & запазване на типизираната абстракция \\
Моделиране, управлявано от отражение & слой на същностите, извличане на метаданни, ергономика на програмния интерфейс & по-малко шаблонен код и по-автоматичен подход, ориентиран към модела & баланс между автоматизация и прозрачност \\
Прецизиране на възможностите & \code{connector\_trait<DB>} и статични проверки & по-точен договор за нови свързващи компоненти & избягване на декоративни labelи за възможности \\
\bottomrule
\end{tabularx}
\end{table}

\subsection{Необходимост от по-строга емпирична оценка}

Една бъдеща версия на работата следва да включва и по-систематичен слой за сравнителен анализ (сравнителен тестing). Той трябва да сравнява не само отделни целеви системи, но и различни режими на изпълнение вътре в библиотеката --- директно изпълнение, \code{prepare()} + повторна употреба, сценарии с масово вмъкване (масово вмъкване) и различни форми на материализация на резултатите. Такъв сравнителен анализ би бил ценен както инженерно, така и научно, защото би показал реалната цена на различните архитектурни избори.

Наред с това, полезно би било да се добави и формализирана оценка на диагностиката на грешките. Една от силните страни на дизайна за ОРС по време на компилация е именно ранното локализиране на грешки. Следователно качеството на съобщенията за грешка и яснотата на пътищата за отказ по време на компилация също са валиден обект на изследване.

\subsection{Обобщение}

Бъдещото развитие на библиотеката не започва от празно място. Напротив, то стъпва върху вече работеща архитектура, която е достатъчно обща, за да поеме нови свързващи компоненти, и достатъчно строга, за да наложи формални ограничения. Именно това прави разработката подходяща както за продължаване като инженерна система, така и за по-нататъшни академични изследвания в областта на моделирането на достъп до данни по време на компилация.

Най-важното заключение от настоящата глава е, че бъдещата работа не трябва да се разбира като поправка на несъстоятелен прототип. Тя е естествено разгръщане на вече защитима архитектурна основа. Това е съществено различие: библиотеката е достигнала ниво, на което разширението е въпрос на дисциплинирано развитие, а не на концептуално преосмисляне.
